% GENERATED FILE. Edit canonical lessons, then run npm run generate:books.
% canonical-source-hash: fnv1a64:3aaa6f5b0879bed6
% canonical-lessons: TA-C107-takeleave, TA-C107-shallwego, TA-C107-seeyoulater, TA-C107-goodjourney, TA-C107-comeagain

\chapter{See You Later}
\label{ch:ta-see-you-later}
\hlchaptermodality{107}

\begin{chapteropening}
\textbf{By the end of this chapter:} \emph{I can take my leave, ask if we should go, say see you later, wish a good journey and ask someone to come again.}

Complete the last lesson of chapter 107: I can take my leave, ask if we should go, say see you later, wish a good journey and ask someone to come again.
\end{chapteropening}

\section[\texorpdfstring{\ta{விடைபெறுகிறேன்} (viṭaipeṟukiṟēṉ)}{விடைபெறுகிறேன் (viṭaipeṟukiṟēṉ)}]{\ta{விடைபெறுகிறேன்} (viṭaipeṟukiṟēṉ) — I take my leave}

\label{lesson:TA-C107-takeleave}



\begin{quote}
Before the new one: say the Tamil for let's go, then the Tamil for a holiday.
\end{quote}

\subsection*{You'll want to know: \ta{விடைபெறுகிறேன்}}
\textbf{\ta{விடைபெறுகிறேன்}} (\emph{viṭaipeṟukiṟēṉ}) — \textquotedblleft{}I take my leave\textquotedblright{}, literally \textquotedblleft{}I receive leave\textquotedblright{}.

\textbf{\ta{விடை}} is \textquotedblleft{}leave, permission\textquotedblright{}. The formal goodbye.

\begin{cousinweb}
\textbf{\ta{விடைபெறுகிறேன்}} (\emph{viṭaipeṟukiṟēṉ}) is native: \emph{viṭai}, \textquotedblleft{}leave, permission to go\textquotedblright{} (DEDR 5393), and \emph{peṟu}, \textquotedblleft{}to get\textquotedblright{} (DEDR 4422): \textquotedblleft{}I get my leave\textquotedblright{}.
\end{cousinweb}

\begin{grammarlens}[title={What you've built}]
A formal way out of a room.
\end{grammarlens}

\subsection*{Your turn}
\begin{itemize}
\raggedright
  \item \emph{Say it:} \emph{viṭaipeṟukiṟēṉ}
  \item \emph{Say it:} \emph{viṭaipeṟukiṟēṉ}, once more
  \item \emph{Say it:} say \emph{viṭaipeṟukiṟēṉ}
  \item \emph{Recall:} say the Tamil for let's go, then the Tamil for a holiday, then say \emph{viṭaipeṟukiṟēṉ} again
\end{itemize}

\begin{tcolorbox}[breakable,colback=teal!4,colframe=teal!35!black,title={Before you move on}]
What does \ta{விடைபெறுகிறேன்} mean? (\textquotedblleft{}I take my leave\textquotedblright{}.) Say it once more.
\end{tcolorbox}
\section[\texorpdfstring{\ta{கிளம்பலாமா}? (kiḷampalāmā?)}{கிளம்பலாமா? (kiḷampalāmā?)}]{\ta{கிளம்பலாமா}? (kiḷampalāmā?) — shall we leave?}

\label{lesson:TA-C107-shallwego}



\begin{quote}
Before the new one: say the Tamil for a holiday, then the Tamil for I take my leave.
\end{quote}

\subsection*{You'll want to know: \ta{கிளம்பலாமா}?}
\textbf{\ta{கிளம்பலாமா}?} (\emph{kiḷampalāmā?}) — \textquotedblleft{}shall we leave?\textquotedblright{}.

\textbf{\ta{கிளம்பு}} is \textquotedblleft{}to set out\textquotedblright{}; \textbf{-\ta{லாமா}} asks \textquotedblleft{}shall we?\textquotedblright{}.

\begin{cousinweb}
\textbf{\ta{கிளம்பலாமா}?} (\emph{kiḷampalāmā?}) is native Dravidian. DEDR 1583 lists Malayalam \emph{kiḷampuka}, \textquotedblleft{}to rise\textquotedblright{}. The entry lists \emph{kiḷampu}, \textquotedblleft{}to rise, to start up\textquotedblright{}: to leave is to get up and go.
\end{cousinweb}

\begin{grammarlens}[title={What you've built}]
A gentle nudge to go.
\end{grammarlens}

\subsection*{Your turn}
\begin{itemize}
\raggedright
  \item \emph{Say it:} \emph{kiḷampalāmā?}
  \item \emph{Say it:} \emph{kiḷampalāmā?}, once more
  \item \emph{Say it:} ask \emph{kiḷampalāmā?}
  \item \emph{Recall:} say the Tamil for a holiday, then the Tamil for I take my leave, then say \emph{kiḷampalāmā?} again
\end{itemize}

\begin{tcolorbox}[breakable,colback=teal!4,colframe=teal!35!black,title={Before you move on}]
What does \ta{கிளம்பலாமா}? mean? (\textquotedblleft{}Shall we leave?\textquotedblright{}.) Say it once more.
\end{tcolorbox}
\section[\texorpdfstring{\ta{அப்புறம்} \ta{பார்க்கலாம்} (appuṟam pārkkalām)}{அப்புறம் பார்க்கலாம் (appuṟam pārkkalām)}]{\ta{அப்புறம்} \ta{பார்க்கலாம்} (appuṟam pārkkalām) — see you later}

\label{lesson:TA-C107-seeyoulater}



\begin{quote}
Before the new one: say the Tamil for I take my leave, then the Tamil for shall we leave.
\end{quote}

\subsection*{You'll want to know: \ta{அப்புறம்} \ta{பார்க்கலாம்}}
\textbf{\ta{அப்புறம்} \ta{பார்க்கலாம்}} (\emph{appuṟam pārkkalām}) — \textquotedblleft{}see you later\textquotedblright{}.

\textbf{\ta{அப்புறம்}} is \textquotedblleft{}afterwards\textquotedblright{}; \textbf{\ta{பார்க்கலாம்}} is \textquotedblleft{}we can see (each other)\textquotedblright{}. To one friend: \textbf{\ta{உன்னை} \ta{அப்புறம்} \ta{பார்க்கிறேன்}} — \textquotedblleft{}I'll see you later\textquotedblright{}, with \textbf{\ta{உன்னை}}, \textquotedblleft{}you\textquotedblright{} as the one it is done to.

\begin{cousinweb}
\textbf{\ta{அப்புறம்} \ta{பார்க்கலாம்}} (\emph{appuṟam pārkkalām}) is native: \emph{appuṟam}, \textquotedblleft{}after that\textquotedblright{}, from the far base \emph{a-} (DEDR 1) and \emph{puṟam}, \textquotedblleft{}side\textquotedblright{} (DEDR 4333), and \emph{pārkkalām}, \textquotedblleft{}let us see\textquotedblright{}, from \emph{pār} (DEDR 4091).
\end{cousinweb}

\begin{grammarlens}[title={What you've built}]
A casual goodbye.
\end{grammarlens}

\subsection*{Your turn}
\begin{itemize}
\raggedright
  \item \emph{Say it:} \emph{appuṟam pārkkalām}
  \item \emph{Say it:} \emph{appuṟam pārkkalām}, once more
  \item \emph{Say it:} say \emph{appuṟam pārkkalām}
  \item \emph{Recall:} say the Tamil for I take my leave, then the Tamil for shall we leave, then say \emph{appuṟam pārkkalām} again
\end{itemize}

\begin{tcolorbox}[breakable,colback=teal!4,colframe=teal!35!black,title={Before you move on}]
What does \ta{அப்புறம்} \ta{பார்க்கலாம்} mean? (\textquotedblleft{}See you later\textquotedblright{}.) Say it once more.
\end{tcolorbox}
\section[\texorpdfstring{\ta{இனிய} \ta{பயணம்} (iṉiya payaṇam)}{இனிய பயணம் (iṉiya payaṇam)}]{\ta{இனிய} \ta{பயணம்} (iṉiya payaṇam) — have a good journey}

\label{lesson:TA-C107-goodjourney}



\begin{quote}
Before the new one: say the Tamil for shall we leave, then the Tamil for see you later.
\end{quote}

\subsection*{You'll want to know: \ta{இனிய} \ta{பயணம்}}
\textbf{\ta{இனிய} \ta{பயணம்}} (\emph{iṉiya payaṇam}) — \textquotedblleft{}have a good journey\textquotedblright{}, literally \textquotedblleft{}a sweet journey\textquotedblright{}.

\textbf{\ta{இனிய}} is \textquotedblleft{}sweet, pleasant\textquotedblright{}.

\textbf{In the family, and next door}

\noindent
\begin{tabularx}{\linewidth}{@{}>{\raggedright\arraybackslash}X>{\raggedright\arraybackslash}X@{}}
\toprule
\textbf{} & \textbf{word} \\
\midrule
Hindi (neighbour) & \dv{शुभ} \dv{यात्रा} (\emph{śubh yātrā}) \\
English & have a good journey \\
\bottomrule
\end{tabularx}

\begin{cousinweb}
\textbf{\ta{இனிய} \ta{பயணம்}} (\emph{iṉiya payaṇam}) is native \emph{iṉiya}, \textquotedblleft{}sweet, pleasant\textquotedblright{} (DEDR 530), and \emph{payaṇam}, \textquotedblleft{}journey\textquotedblright{}, from Sanskrit \emph{prayāṇa} (CDIAL 8732).
\end{cousinweb}

\begin{grammarlens}[title={What you've built}]
A wish for someone setting off.
\end{grammarlens}

\subsection*{Your turn}
\begin{itemize}
\raggedright
  \item \emph{Say it:} \emph{iṉiya payaṇam}
  \item \emph{Say it:} \emph{iṉiya payaṇam}, once more
  \item \emph{Say it:} say \emph{iṉiya payaṇam!}
  \item \emph{Recall:} say the Tamil for shall we leave, then the Tamil for see you later, then say \emph{iṉiya payaṇam} again
\end{itemize}

\begin{tcolorbox}[breakable,colback=teal!4,colframe=teal!35!black,title={Before you move on}]
What does \ta{இனிய} \ta{பயணம்} mean? (\textquotedblleft{}Have a good journey\textquotedblright{}.) Say it once more.
\end{tcolorbox}
\section[\texorpdfstring{\ta{மீண்டும்} \ta{வாருங்கள்} (mīṇṭum vāruṅkaḷ)}{மீண்டும் வாருங்கள் (mīṇṭum vāruṅkaḷ)}]{\ta{மீண்டும்} \ta{வாருங்கள்} (mīṇṭum vāruṅkaḷ) — please come again}

\label{lesson:TA-C107-comeagain}



\begin{quote}
Before the new one: say the Tamil for see you later, then the Tamil for a good journey.
\end{quote}

\subsection*{You'll want to know: \ta{மீண்டும்} \ta{வாருங்கள்}}
\textbf{\ta{மீண்டும்} \ta{வாருங்கள்}} (\emph{mīṇṭum vāruṅkaḷ}) — \textquotedblleft{}please come again\textquotedblright{}.

\textbf{\ta{மீண்டும்}} is \textquotedblleft{}again\textquotedblright{}, with the welcome \textbf{\ta{வாருங்கள்}}.

\begin{cousinweb}
\textbf{\ta{மீண்டும்} \ta{வாருங்கள்}} (\emph{mīṇṭum vāruṅkaḷ}) is native: \emph{mīṇṭum}, \textquotedblleft{}again\textquotedblright{}, from \emph{mīḷ}, \textquotedblleft{}to return\textquotedblright{} (DEDR 4883), and \emph{vā}, \textquotedblleft{}come\textquotedblright{} (DEDR 5270), with the polite ending.
\end{cousinweb}

\begin{grammarlens}[title={What you've built}]
Take leave, shall we go, see you, a good journey, come again.
\end{grammarlens}

\subsection*{Your turn}
\begin{itemize}
\raggedright
  \item \emph{Say it:} \emph{mīṇṭum vāruṅkaḷ}
  \item \emph{Say it:} \emph{mīṇṭum vāruṅkaḷ}, once more
  \item \emph{Say it:} say \emph{mīṇṭum vāruṅkaḷ}
  \item \emph{Recall:} say the Tamil for see you later, then the Tamil for a good journey, then say \emph{mīṇṭum vāruṅkaḷ} again
\end{itemize}

\begin{tcolorbox}[breakable,colback=teal!4,colframe=teal!35!black,title={Before you move on}]
What does \ta{மீண்டும்} \ta{வாருங்கள்} mean? (\textquotedblleft{}Please come again\textquotedblright{}.) Say it once more.
\end{tcolorbox}
